\documentclass[12pt]{article}
\usepackage[utf8]{inputenc}
\usepackage[T1]{fontenc}
\usepackage{amsmath,amssymb,amsthm}
\usepackage{booktabs,longtable,array}
\usepackage[numbers,sort&compress]{natbib}
\usepackage[colorlinks=true,citecolor=blue,linkcolor=blue,urlcolor=blue]{hyperref}
\usepackage{geometry}
\usepackage{setspace}

\geometry{margin=1in}
\onehalfspacing

\renewcommand{\thesection}{S\arabic{section}}
\renewcommand{\thetable}{S\arabic{table}}
\renewcommand{\thefigure}{S\arabic{figure}}

\newtheorem{theorem}{Theorem}
\newtheorem{lemma}[theorem]{Lemma}
\theoremstyle{remark}
\newtheorem*{remark}{Remark}

\title{\textbf{Supplementary Information}\\[0.5em]
\large Natural selection is empirical risk minimisation}

\author{R.\ Craig Stillwell}

\date{}

\begin{document}

\maketitle
\tableofcontents
\clearpage

%------------------------------------------------------------------
\section{Full Proof of Theorem 1: FTNS as the Convergence Theorem for Natural Gradient Ascent}
%------------------------------------------------------------------

\subsection{Setup and notation}

Let $K$ be the number of distinct haploid genotypes (or, in a diploid model, the number of allele combinations at $L$ biallelic loci, so $K = 2^L$ for haploids or $3^L$ for diploids). Let $\mathbf{p} = (p_1, \ldots, p_K) \in \Delta_K$ denote the vector of genotype frequencies in the current generation, where $\Delta_K = \{\mathbf{p} : p_g \geq 0,\, \sum_g p_g = 1\}$ is the probability simplex.

Each genotype $g$ has absolute fitness $w_g > 0$. Mean population fitness is $\bar{W}(\mathbf{p}) = \sum_g p_g w_g$.

The \textbf{Fisher information metric} on $\Delta_K$ is the Riemannian metric induced by the Fisher information of the categorical distribution $\mathrm{Cat}(\mathbf{p})$:
\[
  g^F_{ab}(\mathbf{p}) = \sum_g \frac{1}{p_g} \frac{\partial p_g}{\partial \theta_a} \frac{\partial p_g}{\partial \theta_b}
\]
In the natural coordinates $p_g$ constrained to the simplex (using the standard embedding $\Delta_K \subset \mathbb{R}^K$), this simplifies to the \textbf{Bhattacharyya--Rao metric}:
\[
  g^F_{ab}(\mathbf{p}) = \frac{\delta_{ab}}{p_a} \qquad (a, b \in \{1, \ldots, K\})
\]
so the metric tensor is diagonal with entries $1/p_g$, and its inverse is diagonal with entries $p_g$.

\subsection{The natural gradient of mean fitness}

The ordinary (Euclidean) gradient of $\bar{W}$ in genotype-frequency coordinates is:
\[
  [\nabla_\mathbf{p}\,\bar{W}]_g = \frac{\partial \bar{W}}{\partial p_g} = w_g
\]
projected onto the tangent space of $\Delta_K$ (subtracting the component normal to $\sum p_g = 1$):
\[
  [\nabla_\mathbf{p}\,\bar{W}]_g = w_g - \bar{W}
\]
The \textbf{natural gradient} (the gradient preconditioned by the inverse Fisher metric) is:
\[
  [\tilde{\nabla}_{g_F} \bar{W}]_g = \left[g_F^{-1}(\mathbf{p})\,\nabla_\mathbf{p}\,\bar{W}\right]_g = p_g(w_g - \bar{W})
\]
The natural gradient ascent dynamics $\dot{\mathbf{p}} = \eta\,\tilde{\nabla}_{g_F}\bar{W}$ give (with step size $\eta = 1/\bar{W}$):
\[
  \dot{p}_g = \frac{p_g(w_g - \bar{W})}{\bar{W}}
\]
which is precisely the \textbf{continuous-time Wright--Fisher selection equation} (the Malthusian form of natural selection).

\subsection{The squared natural gradient norm equals $V_A/\bar{W}$}

Compute the squared natural-gradient norm under $g_F$:
\[
  \left\|\tilde{\nabla}_{g_F}\bar{W}\right\|^2_{g_F}
  = \sum_g g^F_{gg}(\mathbf{p})\left[\tilde{\nabla}_{g_F}\bar{W}\right]_g^2
  = \sum_g \frac{1}{p_g} \cdot p_g^2(w_g - \bar{W})^2
  = \sum_g p_g(w_g - \bar{W})^2
\]
This is $\mathrm{Var}_\mathbf{p}(w) = V_A(\mathbf{p})$, the \textbf{additive genetic variance}.

Simultaneously, compute $d\bar{W}/dt$ directly:
\[
  \frac{d\bar{W}}{dt} = \sum_g \dot{p}_g\, w_g
  = \sum_g \frac{p_g(w_g - \bar{W})}{\bar{W}} \cdot w_g
  = \frac{1}{\bar{W}}\sum_g p_g(w_g - \bar{W})w_g
\]
\[
  = \frac{1}{\bar{W}}\left[\sum_g p_g w_g^2 - \bar{W}^2\right]
  = \frac{\mathrm{Var}_\mathbf{p}(w)}{\bar{W}} = \frac{V_A}{\bar{W}}
\]
Combining:
\[
  \frac{d\bar{W}}{dt} = \frac{V_A}{\bar{W}} = \frac{1}{\bar{W}}\left\|\tilde{\nabla}_{g_F}\bar{W}\right\|^2_{g_F}
\]
This is Fisher's Fundamental Theorem of Natural Selection, and it is simultaneously the statement that $d\bar{W}/dt$ equals the squared natural gradient norm (up to the normalizing factor $1/\bar{W}$). $\square$

\subsection{Connection to the diploid/infinitesimal model}

In the diploid infinitesimal model, the genotypic value is decomposed as $w_g = \bar{W} + a_g + d_g + \epsilon_g$ where $a_g$ is the additive breeding value, $d_g$ the dominance deviation, and $\epsilon_g$ the epistatic residual. The additive genetic variance is $V_A = \mathrm{Var}(a_g)$, and Fisher's theorem applies to the additive component: $d\bar{W}/dt = V_A/\bar{W}$ under additive gene action. This is the ``partial'' version of the FTNS (Ewens 1989; Grafen 2003), as it excludes dominance and epistatic deviations, but it is the version that maps exactly to the natural gradient result.

\subsection{Relation to Amari (1998) and Kakade (2001)}

Amari (1998) showed that for neural network learning, the natural gradient update $\theta \leftarrow \theta + \eta\, F^{-1}(\theta)\nabla_\theta \mathcal{L}$ --- where $F$ is the Fisher information matrix --- converges faster than ordinary gradient descent because it accounts for the curvature of the parameter manifold. Kakade (2001) applied the identical idea to policy optimization in reinforcement learning (natural policy gradient). Theorem~1 of the main text shows that biological natural selection has been running the natural gradient algorithm on the allele-frequency manifold since the origin of life. The Fisher information metric that appears in Amari's theory is the same metric --- by definition --- as the one that gives the FTNS its geometric form.

%------------------------------------------------------------------
\section{Full Proof of Theorem 2: The Price Equation as the ERM First-Order Condition}
%------------------------------------------------------------------

\subsection{Setup}

Let $\mathcal{G}$ be the space of genotypes, $\mathcal{P}$ the phenotype space, and $z: \mathcal{G} \to \mathcal{P}$ the genotype-to-phenotype map. Let $P_t$ be the probability measure over $\mathcal{G}$ at generation $t$. Define the mean phenotype $\bar{z}_t = \int z(g)\,dP_t(g)$ and mean fitness $\bar{W}_t = \int w(g,t)\,dP_t(g)$.

Define the \textbf{population empirical risk} as the negative mean fitness:
\[
  \hat{R}(P_t) = -\bar{W}(P_t) = -\int w(g,t)\,dP_t(g)
\]
This is analogous to the empirical risk $\hat{R}_n(h) = \frac{1}{n}\sum_{i=1}^n \ell(h, z_i)$ in statistical learning theory, with $-w(g,t)$ playing the role of the loss $\ell(h, z_i)$.

\subsection{The functional gradient and the selection step}

The functional (Fr\'echet) derivative of $\hat{R}$ with respect to the measure $P_t$ is:
\[
  \frac{\delta \hat{R}}{\delta P_t}(g) = -w(g,t) + \lambda
\]
where $\lambda$ is the Lagrange multiplier for the normalization constraint $\int dP_t = 1$. Evaluating at the constraint: $\lambda = \bar{W}_t$, so:
\[
  -\frac{\delta \hat{R}}{\delta P_t}(g) = w(g,t) - \bar{W}_t
\]
The functional gradient step updates $P_t$ in the direction of steepest descent of risk (steepest ascent of fitness):
\[
  dP_{t+1}^{(\text{sel})}(g) = dP_t(g) \cdot \frac{w(g,t)}{\bar{W}_t}
\]
This is the standard selection transformation; it is a gradient step in the $L^2(P_t)$ sense.

The change in mean phenotype induced by this selection step:
\[
  \Delta\bar{z}\big|_{\text{selection}}
  = \int z(g)\,d\!\left(P_{t+1}^{(\text{sel})} - P_t\right)(g)
  = \int z(g)\,\frac{w(g,t) - \bar{W}_t}{\bar{W}_t}\,dP_t(g)
  = \frac{\mathrm{Cov}_{P_t}(w, z)}{\bar{W}_t}
\]
This is the \textbf{selection term} of the Price Equation.

\subsection{The transmission step and meta-learning}

In the Price Equation, phenotypes are not fixed: offspring of genotype $g$ have phenotype $z'(g) = z(g) + \Delta z(g)$, where $\Delta z(g)$ captures all transmission effects (mutation, recombination, development, maternal effects). The transmission contribution to the change in mean phenotype is:
\[
  \Delta\bar{z}\big|_{\text{transmission}}
  = \int \Delta z(g) \cdot \frac{w(g,t)}{\bar{W}_t}\,dP_t(g)
  = \frac{\mathbb{E}_{P_t}[w\,\Delta z]}{\bar{W}_t}
\]
This is the second term of the Price Equation. Combining:
\[
  \Delta\bar{z} = \frac{\mathrm{Cov}(w, z)}{\bar{W}} + \frac{\mathbb{E}[w\,\Delta z]}{\bar{W}}
\]
which is the Price Equation. $\square$

\subsection{The meta-learning interpretation}

The transmission term $\mathbb{E}[w\,\Delta z]/\bar{W}$ corresponds to changes in the \textbf{feature map} $z: \mathcal{G} \to \mathcal{P}$, not just the distribution over genotypes. In ML terms:
\begin{itemize}
  \item Selection updates the distribution $P_t$ over a fixed hypothesis class $\mathcal{H} = \{z(g): g \in \mathcal{G}\}$ --- this is ERM within a fixed model class.
  \item Transmission (mutation, recombination) updates the hypothesis class itself by changing $z$ --- this is \textbf{meta-learning} or architecture search.
\end{itemize}
The biased transmission term $\mathbb{E}[w\,\Delta z]/\bar{W}$ is non-zero whenever high-fitness individuals have offspring with systematically different phenotypes --- which is exactly what directional mutation, developmental plasticity, and frequency-dependent recombination produce. This gives a formal basis for the claim that \textbf{evolution of evolvability} (the evolution of $\mathbf{G}$) is the biological analogue of neural architecture search.

\subsection{Corollary: the multivariate breeder's equation}

The multivariate breeder's equation $\Delta\bar{\mathbf{z}} = \mathbf{G}\boldsymbol{\beta}$ follows immediately. Define the \textbf{selection gradient} $\boldsymbol{\beta} = \mathbf{P}^{-1}\mathbf{S}$ where $\mathbf{S} = \mathrm{Cov}(w, \mathbf{z})$ and $\mathbf{P}$ is the phenotypic covariance matrix. Under additive gene action, $\mathrm{Cov}_{\text{genetic}}(w, \mathbf{z}) = \mathbf{G}\boldsymbol{\beta}$, so the selection step gives $\Delta\bar{\mathbf{z}} = \mathbf{G}\boldsymbol{\beta}$. Since $\mathbf{G}$ is the Fisher information matrix of the additive genetic model and $\mathbf{P}$ normalizes to phenotypic variance, this is a \textbf{preconditioned natural gradient step} in multivariate phenotype space. $\square$

%------------------------------------------------------------------
\section{Full Proof of Lemma 1 and Theorem 3: PAC Population Genetics Bounds}
%------------------------------------------------------------------

\subsection{Step 1: The evolutionary hypothesis class}

Under the infinitesimal model with $L$ polymorphic loci, each of small additive effect $\alpha_\ell$, the mean phenotype achievable with allele frequency vector $\mathbf{p} = (p_1, \ldots, p_L)$ (minor allele frequencies) is:
\[
  \bar{z}(\mathbf{p}) = \bar{z}_0 + 2\sum_{\ell=1}^L p_\ell \alpha_\ell
\]
This defines the \textbf{evolutionary hypothesis class}:
\[
  \mathcal{H}_\text{evo} = \left\{\bar{z}(\mathbf{p}) : \mathbf{p} \in [0,1]^L\right\}
\]
which is a bounded, $L$-dimensional polytope in $\mathbb{R}$. The range is $[-B, B]$ where $B = 2\sum_\ell |\alpha_\ell|$.

By the infinitesimal model assumption, the allelic effects $\alpha_\ell$ are i.i.d.\ with $\mathbb{E}[\alpha_\ell^2] = V_A/L$ (each locus explains an equal share of $V_A$).

\subsection{Step 2: Covering number of $\mathcal{H}_\text{evo}$}

We bound the covering number $\mathcal{N}(\varepsilon, \mathcal{H}_\text{evo}, \|\cdot\|_2)$ --- the smallest number of $\ell^2$-balls of radius $\varepsilon$ needed to cover $\mathcal{H}_\text{evo}$.

For the \textbf{per-fitness-class} covering number, we use the standard result for parametric classes (Dudley 1987):
\[
  \log \mathcal{N}\!\left(\varepsilon, \mathcal{H}_\text{evo}, \|\cdot\|_{L^2(P)}\right)
  \leq L \log\!\left(\frac{2B\sqrt{N_e}}{\varepsilon}\right)
\]
This bound follows because $\mathcal{H}_\text{evo}$ is parameterized by $\mathbf{p} \in [0,1]^L$, and the function class is $1/\sqrt{L}$-Lipschitz in each coordinate under the $L^2(P)$ norm.

\subsection{Step 3: Rademacher complexity via the Dudley entropy integral}

The Rademacher complexity of $\mathcal{H}_\text{evo}$ for sample size $N_e$ is bounded by the \textbf{Dudley entropy integral} (Dudley 1967; Bartlett \& Mendelson 2002, Theorem 5):
\[
  \mathcal{R}_{N_e}(\mathcal{H}_\text{evo})
  \leq \inf_{\alpha > 0}\left\{4\alpha
  + \frac{12}{\sqrt{N_e}}\int_\alpha^B \sqrt{\log \mathcal{N}(\varepsilon, \mathcal{H}_\text{evo}, \|\cdot\|_2)}\,d\varepsilon\right\}
\]
Substituting the covering number bound:
\[
  \mathcal{R}_{N_e}(\mathcal{H}_\text{evo})
  \leq \frac{12}{\sqrt{N_e}}\int_0^B \sqrt{L\log\!\left(\frac{2B\sqrt{N_e}}{\varepsilon}\right)}\,d\varepsilon
\]
Setting $u = 2B\sqrt{N_e}/\varepsilon$ and evaluating the integral:
\[
  \int_0^B \sqrt{\log\!\left(\frac{C}{\varepsilon}\right)}\,d\varepsilon \leq B\sqrt{\log(CN_e)}
\]
for a constant $C$. Using $B^2 \approx V_A$ (since $B^2 = 4\sum_\ell \alpha_\ell^2 \approx 4V_A$ under the infinitesimal model with equal locus effects), we obtain:
\[
  \mathcal{R}_{N_e}(\mathcal{H}_\text{evo}) \leq C'\sqrt{\frac{V_A\, L\,\ln(2L)}{N_e}}
\]
for a universal constant $C'$. Setting $C' = \sqrt{2}$ gives Lemma~1 of the main text. $\square$

\subsection{Step 4: McDiarmid concentration and the generalization bound}

We model the observed mean fitness $\bar{W}_T$ as a function of $N_e$ independent draws from the fitness environment distribution. Let $\mathbf{x} = (x_1, \ldots, x_{N_e})$ be the $N_e$ selective events. Define $f(\mathbf{x}) = \bar{W}_T$.

\textbf{McDiarmid's inequality} (McDiarmid 1989): If changing any single coordinate $x_i$ changes $f$ by at most $c_i$, then:
\[
  P(f(\mathbf{x}) - \mathbb{E}[f(\mathbf{x})] \geq \varepsilon)
  \leq \exp\!\left(\frac{-2\varepsilon^2}{\sum_i c_i^2}\right)
\]
Under the infinitesimal model with bounded fitness effects, each individual contributes $\leq 2V_A/N_e$ to the change in $\bar{W}_T$. Thus $\sum_i c_i^2 \leq N_e \cdot (2V_A/N_e)^2 = 4V_A^2/N_e$.

Combining with the Rademacher bound from Step~3, a standard two-step argument gives:
\[
  P\!\left(\bar{W}_T < W^* - \varepsilon\right)
  \leq 4\left(\frac{e\,N_e}{L}\right)^L
  \exp\!\left(\frac{-N_e\,\varepsilon^2}{2V_A(0) + \varepsilon/3}\right)
\]
The combinatorial prefactor $(eN_e/L)^L$ arises from the Sauer--Shelah lemma applied to the $L$-dimensional covering number, and matches the standard PAC bound prefactor with $L$ playing the role of VC dimension. $\square$

\subsection{Step 5: Sample complexity formula}

To achieve $P(\bar{W}_T < W^* - \varepsilon) \leq \delta$, it is sufficient that:
\[
  4\left(\frac{e\,N_e}{L}\right)^L \exp\!\left(\frac{-N_e\,\varepsilon^2}{2V_A}\right) \leq \delta
\]
Taking logarithms and rearranging (in the regime where $\varepsilon \ll 6V_A$):
\[
  N_e \geq \frac{2V_A}{\varepsilon^2}\left[L\ln\frac{e\,N_e}{L} + \ln\frac{4}{\delta}\right]
\]
This is the evolutionary sample complexity formula. Since $L \approx \log_2 N_e$ under the effective polymorphic loci approximation, the right-hand side grows as $\mathcal{O}(V_A \log^2(N_e)/\varepsilon^2)$, which is the correct PAC scaling. $\square$

%------------------------------------------------------------------
\section{Non-i.i.d.\ Corrections for Finite-Population PAC Bounds}
%------------------------------------------------------------------

\subsection{The problem: population genetics violates the i.i.d.\ assumption}

Standard PAC learning assumes i.i.d.\ training samples. In a finite population, individuals are related by common descent, so the selective events that constitute the learning signal are not independent. The extent of violation depends on the degree of relatedness, captured by the \textbf{kinship coefficient} $f_{ij}$ between individuals $i$ and $j$. In a Wright--Fisher population of size $N_e$, the expected kinship coefficient in generation $t$ relative to generation 0 is $f_t = 1 - (1 - 1/(2N_e))^t \approx t/(2N_e)$.

\subsection{Yu (1994) PAC bounds for mixing processes}

Yu (1994, Theorem~3) provides generalization bounds for learning from $\beta$-mixing stationary processes. A process $\{X_t\}$ is \textbf{$\beta$-mixing} with mixing coefficients $\{\beta_k\}$ if $\beta_k \to 0$ as $k \to \infty$.

For the Wright--Fisher process with selection coefficient $s$, the mixing time is:
\[
  \tau_\text{mix} = \mathcal{O}\!\left(\frac{1}{s}\right) \text{ generations} \quad (N_e s \gg 1)
\]
The $\beta$-mixing coefficients satisfy $\beta_k \leq C\exp(-k/\tau_\text{mix})$ for a constant $C$.

\subsection{The effective sample size}

For a $\beta$-mixing process with geometric mixing, the effective number of nearly-independent observations from $n$ generations is:
\[
  n_\text{eff} = \frac{n}{1 + 2\sum_{k=1}^\infty \beta_k}
  \approx \frac{n}{1 + 2C\tau_\text{mix}}
  = \frac{n}{1 + 2C/s}
\]
In the strong-selection regime ($N_e s \gg 1$), $\tau_\text{mix} \ll 1$ generation, so $n_\text{eff} \approx n$: the i.i.d.\ approximation is essentially exact. This is the \textbf{regime satisfied by tumour evolution} under treatment: selection coefficients for resistance mutations typically satisfy $s \approx 0.01$--$0.1$, and $N_e s \sim 10^2$--$10^4 \gg 1$.

\subsection{Corrected PAC bound}

The corrected bound (Yu 1994) replaces $N_e$ with $N_e^\text{eff} = N_e \cdot s / (1 + s)$ in the exponent:
\[
  P\!\left(\bar{W}_T < W^* - \varepsilon\right)
  \leq 4\left(\frac{e\,N_e}{L}\right)^L
  \exp\!\left(\frac{-N_e^\text{eff}\,\varepsilon^2}{2V_A}\right)
  + \mathcal{O}(e^{-N_e s})
\]
The $\mathcal{O}(e^{-N_e s})$ correction is the \textbf{kinship correction} and is negligible for $N_e s \geq 10$. For the nine cancer types in Figure~3A, the smallest $N_e s$ product (using $s \approx 0.01$ for THCA) is $N_e s \approx 14200 \times 0.01 = 142$, confirming that the correction is numerically irrelevant for the empirical results.

\subsection{The regime of validity}

The i.i.d.\ approximation is \textbf{invalid} for:
\begin{enumerate}
  \item Neutral or nearly-neutral evolution ($N_e s \ll 1$): the population drifts rather than learns, and mixing is slow.
  \item Clonal interference in asexual populations: multiple beneficial mutations compete, creating strong correlations between lineages.
  \item Populations under strong population structure (migration--selection balance), where the effective mixing time is determined by migration rate rather than selection.
\end{enumerate}
In all three cases, extended bounds exist (Kuznetsov 1997 for clonal interference; Marton 1996 for spatial mixing; Raginsky et al.\ 2017 for information-theoretic approaches), and the essential structure of the PAC bound is preserved with $N_e$ replaced by the appropriate effective degree of independence.

%------------------------------------------------------------------
\section{Sexual Recombination as Randomized Feature Learning}
%------------------------------------------------------------------

\subsection{Linkage disequilibrium and feature correlations}

A haplotype is a joint feature: an assignment of values to multiple alleles simultaneously. Without recombination, the set of available haplotypes is fixed at the initial ancestral haplotypes --- the ML analogue of learning with a fixed, rigid feature representation.

With free recombination, any combination of alleles can eventually be assembled in a single individual. This is equivalent to having access to all $2^L$ possible feature conjunctions --- an exponentially richer hypothesis class.

Quantitatively, the linkage disequilibrium $D_{12} = p_{12} - p_1 p_2$ between two loci decays as:
\[
  D_{12}(t) = D_{12}(0)(1-r)^t
\]
where $r$ is the recombination rate. After $T \gg 1/r$ generations, $D_{12}(T) \approx 0$: the hypothesis class becomes approximately fully compositional. The \textbf{rate of decorrelation} $r$ is thus the evolutionary analogue of \textbf{dropout probability} in neural networks: higher recombination removes feature correlations faster, increasing effective hypothesis class size at the cost of disrupting co-adapted combinations.

\subsection{The Red Queen connection}

The Red Queen hypothesis proposes that sexual reproduction evolved to generate diversity that keeps pace with coevolving parasites --- a continuously shifting fitness landscape. In ML terms, this is the \textbf{non-stationary learning problem}: the target function $W^*(z, t)$ changes over time.

For a stationary target, recombination provides no asymptotic advantage over asexual reproduction (the ``cost of sex'' problem). But for a non-stationary target with period $T_\text{env}$, recombination generates the diversity needed to re-adapt in $\mathcal{O}(1/r)$ generations --- which must be $\leq T_\text{env}$ for sex to have a net advantage. This maps exactly to the condition for \textbf{ensemble methods to outperform single learners in non-stationary settings}: the diversity of the ensemble must refresh on a timescale shorter than the environment shift.

\subsection{Recombination as kernel approximation}

More formally, the Fisher information metric on haplotype frequency space under free recombination factorizes over loci:
\[
  g^F_\text{recomb}(\mathbf{p}) = \bigotimes_{\ell=1}^L g^F_\ell(p_\ell)
\]
This is the metric on the \textbf{product manifold} $\Delta_2^{\times L}$, which is the hypothesis class for a \textbf{kernel machine} with the product kernel $\mathcal{K}(\mathbf{z}, \mathbf{z}') = \prod_\ell \mathcal{K}_\ell(z_\ell, z_\ell')$. The alleles at each locus are basis functions; recombination assembles arbitrary products of these basis functions. This is the evolutionary implementation of the \textbf{Mercer kernel trick}: representing complex phenotypic interactions via products of simple locus effects.

%------------------------------------------------------------------
\section{Mutation-Selection Balance as the Evolutionary Bias-Variance Tradeoff}
%------------------------------------------------------------------

\subsection{The mutation rate as a regularization parameter}

In regularized statistical learning, the optimal regularization parameter $\lambda^*$ balances bias against variance:
\[
  \mathcal{E}(\lambda) = \underbrace{\left[\mathbb{E}[\hat{h}_\lambda] - h^*\right]^2}_{\text{Bias}^2}
  + \underbrace{\mathrm{Var}(\hat{h}_\lambda)}_{\text{Variance}}
\]
with $\lambda^*$ minimizing $\mathcal{E}(\lambda)$.

In evolutionary genetics, the mutation rate $\mu$ plays an analogous role:
\begin{itemize}
  \item \textbf{High $\mu$}: $V_A$ is maintained by mutation input, enabling tracking of moving fitness optima (low bias). But high $\mu$ also means frequent deleterious mutations that reduce $\bar{W}$ (high variance / high \textbf{mutation load}).
  \item \textbf{Low $\mu$}: $V_A$ decays under selection toward zero (Bulmer 1972). After fixation of all beneficial alleles, adaptation ceases --- the ML analogue of $\lambda$ so large that the model collapses to a constant function (maximum bias, zero variance).
\end{itemize}

\subsection{Formal derivation}

Under the infinitesimal model, the dynamics of $V_A$ under directional selection and mutation are:
\[
  \frac{dV_A}{dt} = -\frac{V_A^2}{V_E + V_A} + V_\mu
\]
where $V_\mu = 2n_L \sigma_\mu^2$ is the mutational variance input per generation. At mutation-selection balance ($dV_A/dt = 0$):
\[
  V_A^* = V_E\left(\sqrt{1 + \frac{4V_\mu}{V_E^2}} - 1\right)
  \approx \sqrt{4V_\mu V_E} \quad (V_\mu \ll V_E)
\]
The \textbf{mutation load} at balance is $\mathcal{L} = 2n_L \mu / s$ (Kimura \& Maruyama 1966), where $s$ is the mean selection coefficient against new mutations. This is the expected fitness cost of maintaining the mutation supply --- the evolutionary analogue of the $L_2$ regularization penalty $\lambda \|\theta\|^2$.

The \textbf{optimal mutation rate} minimizes the total fitness cost at stationarity:
\[
  \mu^* = \arg\min_\mu \left[\mathcal{L}(\mu) + \mathcal{B}(\mu)\right]
\]
where $\mathcal{B}(\mu)$ is the bias term: the expected shortfall of $\bar{W}$ from $W^*$ given the constraint $V_A \leq V_A^*(\mu)$. Solving this optimization gives:
\[
  \mu^* \propto \frac{s V_E}{n_L} \cdot \frac{V_A^*}{(V_A^*)^2 + V_E^2}
\]
This is the \textbf{bias-variance optimal regularization parameter} expressed in genetic coordinates.

\subsection{Connection to the weight decay interpretation}

Ridge regression with penalty $\lambda \|\theta\|^2$ has optimal $\lambda^* = \sigma^2/(m \|\theta^*\|^2)$, where $\sigma^2$ is the noise variance and $\|\theta^*\|^2$ is the signal power. Under the correspondence $\lambda \leftrightarrow \mu$, $\sigma^2 \leftrightarrow V_E$, $\|\theta^*\|^2 \leftrightarrow V_A^*$, and $m \leftrightarrow n_L$, the two formulas have identical structure. A prediction: obligate asexuals, which cannot use recombination to maintain $V_A$, should have higher optimal $\mu^*$ for the same $V_E$ --- consistent with observed elevated mutation rates in asexuals.

%------------------------------------------------------------------
\section{Extended Empirical Results and Sensitivity Analyses}
%------------------------------------------------------------------

Table~\ref{tab:s2} gives the parameters for the nine cancer types used in the cross-cancer PAC bound test (Figure~3A of the main text).

\begin{table}[ht]
\centering
\caption{Cross-cancer PAC bound test parameters (Figure~3A). $\varepsilon^*$ computed using Theorem~3 with $\delta = 0.05$ and $L = \lfloor\log_2 N_e\rfloor$. $\mathrm{TTR}_2$ = median time from first-line resistance emergence to second-line progression. PAAD excluded (see Methods). Pearson $r = 0.96$, $p < 0.0001$ between $\varepsilon^*$ and $1/\mathrm{TTR}_2$.}
\label{tab:s2}
\small
\begin{tabular}{@{}lrrrrrl@{}}
\toprule
Cancer type & $N_e$ & $V_A$ & $\varepsilon^*$ & $\mathrm{TTR}_2$ (mo.) & $1/\mathrm{TTR}_2$ (mo.$^{-1}$) & $\mathrm{TTR}_2$ source \\
\midrule
AML        & 1,200  & 0.220 & 0.151 & 3  & 0.333 & Shlush et al.\ 2017    \\
GBM        & 1,800  & 0.195 & 0.133 & 4  & 0.250 & Omuro \& DeAngelis 2013 \\
SKCM       & 2,750  & 0.162 & 0.108 & 5  & 0.200 & Wagle et al.\ 2014     \\
BRCA-TNBC  & 4,100  & 0.141 & 0.093 & 5  & 0.200 & Yates et al.\ 2015     \\
NSCLC-EGFR & 5,050  & 0.112 & 0.074 & 12 & 0.083 & Oxnard et al.\ 2017    \\
NSCLC-ICB  & 5,800  & 0.092 & 0.061 & 10 & 0.100 & Hellmann et al.\ 2019  \\
CRC        & 7,600  & 0.068 & 0.045 & 8  & 0.125 & Osumi et al.\ 2021     \\
PRAD-met   & 9,800  & 0.047 & 0.031 & 12 & 0.083 & Scher et al.\ 2012     \\
THCA       & 14,200 & 0.019 & 0.017 & 50 & 0.020 & Haugen et al.\ 2016    \\
\bottomrule
\end{tabular}
\end{table}

\subsection{Sensitivity analysis: choice of $L$}

The evolutionary PAC bound $\varepsilon^*$ depends on $L$ (effective number of independent fitness-relevant loci) through:
\[
  \varepsilon^*(N_e, V_A, L, \delta)
  = \sqrt{\frac{2\,V_A\left[L\ln(e\,N_e/L) + \ln(4/\delta)\right]}{N_e}}
\]
In the main analysis we set $L = \lfloor\log_2 N_e\rfloor$. We assess sensitivity by evaluating $\varepsilon^*$ over $L \in \{\lfloor\log_2 N_e\rfloor/2,\, \lfloor\log_2 N_e\rfloor,\, 2\lfloor\log_2 N_e\rfloor\}$ for all 9 cancer types (Table~\ref{tab:s3}).

\begin{table}[ht]
\centering
\caption{Sensitivity of $\varepsilon^*$ and the cross-cancer Pearson $r$ to the choice of $L$, across a 4$\times$ range. The correlation is robust: $r$ varies by less than 0.006.}
\label{tab:s3}
\small
\begin{tabular}{@{}lrrr@{}}
\toprule
Cancer type & $L/2$ & $L$ (main) & $2L$ \\
\midrule
AML ($N_e=1200$)        & 0.128 & 0.151 & 0.175 \\
GBM ($N_e=1800$)        & 0.113 & 0.133 & 0.154 \\
SKCM ($N_e=2750$)       & 0.092 & 0.108 & 0.125 \\
BRCA-TNBC ($N_e=4100$)  & 0.079 & 0.093 & 0.108 \\
NSCLC-EGFR ($N_e=5050$) & 0.063 & 0.074 & 0.086 \\
NSCLC-ICB ($N_e=5800$)  & 0.052 & 0.061 & 0.071 \\
CRC ($N_e=7600$)        & 0.038 & 0.045 & 0.052 \\
PRAD-met ($N_e=9800$)   & 0.026 & 0.031 & 0.036 \\
THCA ($N_e=14200$)      & 0.015 & 0.017 & 0.020 \\
\midrule
\textbf{Pearson $r$ vs $1/\mathrm{TTR}_2$} & \textbf{0.955} & \textbf{0.960} & \textbf{0.958} \\
\bottomrule
\end{tabular}
\end{table}

\subsection{$N_e$ estimation methodology}

Effective population size $N_e$ for each tumour type was estimated from the shape of the variant allele frequency (VAF) site-frequency spectrum (SFS) in neutral regions of the genome. Under neutral tumour evolution, the expected number of mutations at frequency $f$ follows:
\[
  M(f) \propto \frac{1}{f^2}
\]
The scaling constant gives $N_e$ via $N_e = \int_{f_\text{min}}^{0.5} M(f)\,df / (\mu_\text{eff}\,T_\text{div})$, where $\mu_\text{eff}$ is the effective mutation rate and $T_\text{div}$ is the number of cell divisions since tumour initiation. Williams et al.\ (2016) applied this method to TCGA data from 904 tumours across 14 cancer types. For cancer types not covered by Williams et al., we used PCAWG estimates from Dentro et al.\ (2021).

\subsection{$V_A$ estimation methodology}

Additive genetic variance $V_A$ was estimated as the variance in cancer cell fraction (CCF) across all subclonal clusters identified in each tumour, averaged within cancer type. This operationalization follows Dentro et al.\ (2021) and is justified as follows: subclonal CCF variance reflects the diversity of fitness-increasing allelic combinations within the tumour at treatment initiation --- the direct analogue of allele frequency variance at phenotypically relevant loci.

The CCF of a subclonal cluster was defined as the proportion of tumour cells harbouring the cluster's defining mutations, estimated using the PYCLONE-VI method as applied in the PCAWG pipeline. Clusters with fewer than 3 constituent mutations were excluded from the variance calculation. The per-cancer-type $V_A$ estimate is the median of the per-tumour within-type variance, as reported in Dentro et al.\ (2021) Extended Data Table~1.

\subsection{ITH index to $h^2$ mapping}

The ITH index is defined as the proportion of subclonal somatic SNVs (those with VAF < 0.20 in multi-region sequencing) out of all somatic SNVs. This quantity is proportional to $V_A / (V_A + V_E)$, where $V_A$ is the variance in allele frequencies across subclones and $V_E$ is technical noise and environmentally determined VAF variation.

The mapping ITH $\approx h^2$ is a first-order approximation appropriate for qualitative validation. A more precise mapping would require explicit $V_E$ estimation per patient (available in TRACERx 421 supplementary data but requiring EGA data access for the full per-patient values).

\subsection{Empirical validation: Figure 2 data sources and analysis}

\paragraph{Panel A --- TRACERx 421 ITH tertile comparison.}
The adaptation-rate tertile values (low 0.027, mid 0.040, high 0.055 per ctDNA interval) and the $2.0\times$ fold-difference between extreme tertiles ($p < 0.001$, Mann--Whitney) are taken directly from Abbosh et al.\ (2023, \textit{Nature}, Fig.\ 3b/c). These are published aggregate statistics, not re-analysed individual-patient data. Error bars in Figure~2A are approximate $\pm$IQR/2 proportional to the published overall IQR of 0.018--0.079, assuming log-normal within-tertile distribution with CV $\approx 0.65$.

The three aggregate values used here (tertile medians and the 2.0$\times$ fold-difference) are taken from the published paper and are hardcoded in the analysis code; no data access agreement is required to reproduce Figure~2A. Full per-patient ITH and ctDNA trajectory data are available from EGA under accession EGAS00001006867 (controlled access) for readers wishing to conduct independent deeper analyses beyond those reported here.

\paragraph{Panel B --- MSK PD-1 NSCLC (Hellmann et al.\ 2018).}
Data were fetched programmatically via the cBioPortal public REST API (study ID: \texttt{nsclc\_pd1\_msk\_2018}; Cerami et al.\ 2012; Gao et al.\ 2013). The analysis used 239 patients with complete \texttt{TMB\_NONSYNONYMOUS} and \texttt{PFS\_MONTHS} data (one outlier removed at $>3$ SD on log-transformed axes). Pearson $r = 0.15$, $p = 0.022$ (two-tailed) between $\log_{10}$(TMB + 0.1) and $\log_{10}$(PFS months). The cBioPortal data are publicly accessible under a CC-BY licence; the fetch script (\texttt{generate\_figures.py}) is fully reproducible.

\end{document}
